\chapterbackmatter{Solutions to Weinberg Exercises}
\ExerciseEditorialNote

\begin{enumerate}
\WeinbergSolution{1}
In \(D\) space-time dimensions a canonically normalized real scalar has
\[
 [\phi]=\frac{D-2}{2},\qquad [\partial_\mu]=1.
\]
An operator is power-counting renormalizable when its dimension is at most
\(D\); it is superrenormalizable when the inequality is strict.

In \(D=3\), \([\phi]=1/2\).  Apart from a constant and the kinetic term, the
complete polynomial list is
\[
 \phi,\quad \phi^2,\quad\phi^3,\quad\phi^4,\quad\phi^5,\quad\phi^6.
\]
The first five monomials are superrenormalizable and \(\phi^6\) is
marginal.  No additional derivative interaction is allowed: after
integration by parts the first candidate is
\(\phi^n(\partial\phi)^2\), of dimension \(3+n/2\), and only \(n=0\)
gives the kinetic term.

In \(D=6\), \([\phi]=2\), so the list is
\[
 \phi,\qquad \phi^2,\qquad \phi^3,
\]
together with the kinetic term and a constant.  The cubic interaction is
marginal; the lower powers are superrenormalizable.

The two-dimensional case is exceptional because \([\phi]=0\).  Every
polynomial potential is superrenormalizable:
\[
 V(\phi)=\sum_{n\geq0}g_n\phi^n,\qquad [g_n]=2.
\]
At the level of local power counting one may also write
\[
 \frac12 Z(\phi)(\partial_\mu\phi)(\partial^\mu\phi),
\]
with an arbitrary polynomial \(Z\); all of its terms are marginal.  For one
scalar and nonsingular \(Z\), however, the field redefinition
\(\dd\chi/\dd\phi=\sqrt{Z(\phi)}\) restores a canonical kinetic term.  Thus an
operator basis modulo field redefinitions consists simply of the kinetic
term and the arbitrary potential.  Terms containing more than two
derivatives are nonrenormalizable in \(D=2\).

\WeinbergSolution{2}
Let \(G\) be the two-loop electron self-energy graph in which the internal
photon joins two points of the electron line.  There are two logarithmically
divergent one-loop vertex subgraphs, \(\gamma_1\) and \(\gamma_2\), according
to which external electron segment is included.  They share an internal
electron line and therefore overlap.  In particular, a BPHZ forest may
contain \(\gamma_1\) or \(\gamma_2\), but not both.

For a vertex subgraph let \(t_\gamma\) denote its degree-zero Taylor
operation in the momenta entering the subgraph.  If no other subdivergence is
present, the subdivergence-subtracted integrand is
\[
 \bar R I_G=I_G-t_{\gamma_1}I_G-t_{\gamma_2}I_G.
\]
Each subtraction replaces the corresponding one-loop vertex by its local
divergent part,
\[
 t_{\gamma_i}\Gamma^\mu_{\text{1 loop}}
   =L_{\rm \dd iv}\gamma^\mu .
\]
This is exactly the graph produced by inserting the one-loop QED vertex
counterterm at that location.  The Ward identity gives
\(L_{\rm \dd iv}=(Z_2-1)_{\rm \dd iv}\), so the same coefficient is already fixed
by electron field renormalization.

Consider a subintegration in which the loop momentum of \(\gamma_1\) becomes
large while the other independent momentum is held fixed.  The difference
\((1-t_{\gamma_1})I_G\) falls by one additional inverse power and is
convergent.  The analogous statement holds for \(\gamma_2\).  There is no
double subtraction \(t_{\gamma_1}t_{\gamma_2}I_G\), precisely because the
two subgraphs overlap rather than form a forest.  Generic linear
combinations of the two loop momenta are already superficially convergent.

After these two local vertex subtractions, only the overall electron
self-energy divergence remains.  Its degree is one, so Lorentz invariance
restricts its Taylor polynomial to
\[
 t_G\bar R I_G=A_{\rm \dd iv}+B_{\rm \dd iv}\,\ii\sl p.
\]
The mass and field-strength counterterms cancel these two terms.  Equivalently,
the full forest formula is
\[
 \boxed{\quad
 R I_G=(1-t_G)
 \bigl(I_G-t_{\gamma_1}I_G-t_{\gamma_2}I_G\bigr),
 \quad}
\]
with each Taylor term represented by an ordinary QED counterterm graph.
Thus every ultraviolet subintegration, including the two overlapping
regions, and the overall integration are finite.

\WeinbergSolution{3}
It is useful to make both the polynomial and the convergence of the
remainder explicit.  Work in Euclidean momentum and define
\(\Pi_E(q)\) as the additive correction to the inverse scalar propagator.
The closed fermion loop gives
\[
 \Pi_E^{(1)}(q)
 =4g^2\int_0^1\dd x\int\frac{\dd^4\ell}{(2\pi)^4}
 \frac{m_\psi^2-\ell^2+x(1-x)q^2}
 {\bigl[\ell^2+m_\psi^2+x(1-x)q^2\bigr]^2},
\]
where the shift \(\ell=k+xq\) has been made.  A regulator is understood in
the two terms that follow.

Put \(a=x(1-x)\), \(u=\ell^2+m_\psi^2\), and
\[
 F(\ell,q^2;x)=\frac{m_\psi^2-\ell^2+a q^2}{(u+a q^2)^2}.
\]
Taylor's formula through first order gives
\[
 \Pi_E^{(1)}(q)=A+Bq^2+\Pi_{\rm conv}(q^2),
\]
where
\[
 A=4g^2\int_0^1\dd x\int\frac{\dd^4\ell}{(2\pi)^4}
       \frac{m_\psi^2-\ell^2}{u^2},
\]
\[
 B=4g^2\int_0^1\dd x\int\frac{\dd^4\ell}{(2\pi)^4}
       \frac{a(3\ell^2-m_\psi^2)}{u^3},
\]
and
\[
 \boxed{\;
 \Pi_{\rm conv}(q^2)=4g^2\int_0^1\dd x\int\frac{\dd^4\ell}{(2\pi)^4}
 \left[
 F(\ell,q^2;x)-\frac{m_\psi^2-\ell^2}{u^2}
 -q^2\frac{a(3\ell^2-m_\psi^2)}{u^3}
 \right].\;}
\]
The first line in the bracket has had its constant and \(q^2\) Taylor
coefficients removed.  Consequently the bracket is
\(O(q^4/\ell^6)\) for large \(\ell\), and the four-dimensional integral is
absolutely ultraviolet convergent.

The divergent constants \(A\) and \(B\) are local: they multiply
\(\phi^2\) and \((\partial\phi)^2\), respectively.  A scalar mass
counterterm and a scalar field-strength counterterm therefore remove them.
Analytic continuation \(q_E^2=-q^2-\ii0\) gives the requested Minkowski
self-energy and its physical cut.

\WeinbergSolution{4}
Choose an operator basis modulo integrations by parts and the leading QED
equations of motion.  The first corrections have dimension five.  For one
Dirac field there are two independent Hermitian dipoles,
\[
 \boxed{\;
 \mathcal L^{(5)}
 =\frac{c_M}{M}\bar\psi\sigma^{\mu\nu}\psi F_{\mu\nu}
 +\frac{c_E}{M}\ii\bar\psi\sigma^{\mu\nu}\gamma_5\psi F_{\mu\nu}.
 \;}
\]
The first is the magnetic Pauli operator and the second is the electric
dipole operator.  The latter would be forbidden by P or T, but those
symmetries were not assumed.  Operators such as
\(\bar\psi D^2\psi\) reduce, by
\[
 (\ii\slashed D)^2=-D^2-\frac{e}{2}\sigma^{\mu\nu}F_{\mu\nu},
\]
to these dipoles and lower-dimensional terms.

At dimension six an economical on-shell basis is most transparent in
chiral notation:
\[
\boxed{\;
\begin{aligned}
 \mathcal L^{(6)}=\frac1{M^2}\bigl[&
 c_{LL}(\bar\psi_L\gamma_\mu\psi_L)
        (\bar\psi_L\gamma^\mu\psi_L)\\
 {}+{}&
 c_{RR}(\bar\psi_R\gamma_\mu\psi_R)
        (\bar\psi_R\gamma^\mu\psi_R)\\
 {}+{}&
 c_{LR}(\bar\psi_L\gamma_\mu\psi_L)
        (\bar\psi_R\gamma^\mu\psi_R)
 \bigr]
\end{aligned}\;}
\]
Fierz identities reduce all other four-fermion contractions for a single
species to these three.  Unequal \(c_{LL}\) and \(c_{RR}\) and the equivalent
vector--axial form make clear that discrete symmetries need not be present.

An off-shell basis may instead retain the two anapole/charge-radius
operators
\[
 (\bar\psi\gamma_\mu\psi)\partial_\nu F^{\nu\mu},
 \qquad
 (\bar\psi\gamma_\mu\gamma_5\psi)\partial_\nu F^{\nu\mu},
\]
and the photon operator
\((\partial_\mu F^{\mu\nu})(\partial^\rho F_{\rho\nu})\).
The Maxwell equation
\(\partial_\nu F^{\nu\mu}=e\bar\psi\gamma^\mu\psi\) converts them into
the four-fermion basis above, so they must not be counted twice.  A cubic
Abelian field-strength scalar vanishes in four dimensions, while the
remaining two-fermion derivative terms reduce by integration by parts and
the Dirac equation.  Hence the two dipoles are the complete leading
\(1/M\) set, and the three current--current structures give a complete
next-to-leading \(1/M^2\) on-shell set.
\end{enumerate}
